\begin{abstract}
Fake-news sources can imitate established media without changing how often they publish false content. We study this mechanism with FAKENEWS-ABM, an endorsement-based agent model in which users learn about sources, select one source per period, and exchange word-of-mouth (WOM) recommendations. Across 2,903 runs, we separate source popularity from actual false reposting and examine copy scope, intervention timing, memory, truth-sensitive WOM, contacts, and reach. The baseline reproduces configured false-publication probabilities. Penalising false recommendations while rewarding true ones reduces late-run false reposts by 5.26 percentage points relative to discovery-only WOM. Full attribute copying produces the largest target-selection gain (8.54 points) but reduces false reposting because credibility also controls objective false-publication probability. Early copying of all non-credibility attributes instead preserves the target's false-news propensity. A 30-seed confirmation finds a 2.89-point increase in false reposts under memory 25, false-news penalisation, and no reward for true news. Effects weaken with delayed intervention or short memory. Common seeds reproduce condition rankings and usually narrow uncertainty. The results show why popularity is an inadequate misinformation proxy and identify memory, timing, and credibility coupling as central modelling considerations.
\end{abstract}

\begin{keywords}
agent-based simulation; fake news; misinformation; word of mouth; endorsements; memory; scenario analysis
\end{keywords}
